\documentclass[10pt]{article}
\usepackage[UTF8]{ctex}
\usepackage{amsmath, amssymb}
\usepackage{geometry}
\geometry{a4paper, margin=1.5cm}
\usepackage{enumitem}

\title{导数入门：从变化率到微积分基石}
\author{}
\date{}

\begin{document}

\maketitle

\begin{center}
\textit{——写给高二文科生的学习笔记}
\end{center}

亲爱的同学，你好！导数听起来高深，其实它就是我们身边描述“变化快慢”的数学工具。比如汽车速度计显示的就是位置随时间的变化率；爬山时坡度陡不陡，就是高度随水平距离的变化率。导数就是用来精确计算这些“瞬间变化率”的。我们将从最基础的例子开始，一步步走进导数的世界，你会发现它并没有想象中那么神秘。

\section*{第一部分：预备知识——极限与无穷小}

在正式学习导数之前，我们需要先理解两个关键概念：\textbf{极限}和\textbf{无穷小量}。它们是理解导数的基础。

\subsection*{1.1 极限——无限接近但不等}

\textbf{直观理解}：当你慢慢走向一扇门，你与门的距离越来越小，最后几乎为 0，但还没碰到门的那一瞬间，这个距离就是“趋近于 0”。数学上，我们关心当自变量 $x$ 无限接近某个数 $x_0$ 时，函数 $f(x)$ 会无限接近一个确定的数 $A$，这个 $A$ 就是极限。

\textbf{为什么要引入极限？} 因为有些函数在 $x_0$ 处可能没有定义（比如分母为 0），但我们想知道它附近的变化趋势，极限就能帮我们描述这种“趋势”。

例如，考虑函数 $f(x)=2x+1$，当 $x$ 越来越接近 2 时，$f(x)$ 越来越接近 5。我们可以写为：
\[
\lim_{x \to 2} (2x+1) = 5
\]
这里 $x=2$ 时函数有定义且等于 5，所以极限就等于函数值。

\textbf{计算极限的基本方法}：
\begin{enumerate}
    \item \textbf{直接代入}：如果函数在 $x_0$ 处有定义且图像连续（比如多项式），那么极限就等于函数值。如上例。
    \item \textbf{消去零因子}：当分子分母同时趋近于 0 时，需要先化简。比如：
    \[
    \lim_{x \to 1} \frac{x^2-1}{x-1}
    \]
    当 $x=1$ 时，分母为 0，不能直接代入。但在极限过程中，$x$ 是无限接近 1 但永远不等于 1 的，所以 $x-1 \neq 0$，我们可以安全地约去分子分母的公因子 $x-1$：
    \[
    \frac{x^2-1}{x-1} = \frac{(x-1)(x+1)}{x-1} = x+1 \quad (x \neq 1)
    \]
    这时，当 $x \to 1$ 时，$x+1$ 趋近于 2。所以极限为 2。
\end{enumerate}

\textbf{直观理解：高阶小量可以忽略}

在极限过程中，我们经常遇到像 $\Delta x$ 这样趋于 0 的量，它被称为无穷小量。而 $(\Delta x)^2$ 比 $\Delta x$ 更快地趋于 0（比如 $\Delta x=0.1$ 时，$(\Delta x)^2=0.01$，小了 10 倍），我们称 $(\Delta x)^2$ 是比 $\Delta x$ 更高阶的无穷小。

在求极限时，高阶小量常常可以忽略不计，因为它们在最终结果中不贡献任何值。例如，在 $(x+\Delta x)^2 = x^2 + 2x\Delta x + (\Delta x)^2$ 中，当 $\Delta x \to 0$ 时，$2x\Delta x$ 是主要部分，而 $(\Delta x)^2$ 可以忽略。这正是后面求导时只保留线性项的依据。

\subsection*{1.2 无穷小量——微小的变化}

我们用 $\Delta x$ 表示自变量一个微小的改变量（读作“delta x”）。它不是一个固定的数，而是代表一个可以任意小的变化。比如从 $x=3$ 变到 $3.001$，那么 $\Delta x = 0.001$。在极限过程中，我们让 $\Delta x$ 趋于 0，但永远不为 0。这个思想贯穿导数始终。

\subsection*{1.3 极限的四则运算法则（非常重要，后面会反复用到）}

\begin{itemize}
    \item 和的极限：$\lim [f(x)+g(x)] = \lim f(x) + \lim g(x)$
    \item 差的极限：$\lim [f(x)-g(x)] = \lim f(x) - \lim g(x)$
    \item 积的极限：$\lim [f(x)g(x)] = \lim f(x) \cdot \lim g(x)$
    \item 商的极限：$\lim \frac{f(x)}{g(x)} = \frac{\lim f(x)}{\lim g(x)}$（分母极限不为 0）
    \item 常数因子可提出：$\lim [c \cdot f(x)] = c \cdot \lim f(x)$
\end{itemize}

这些法则让我们在求极限时可以分步计算，非常方便。

\section*{第二部分：导数的定义——从平均变化率到瞬时变化率}

\subsection*{2.1 一个熟悉的例子：平均速度与瞬时速度}

想象一辆车在 $t$ 时刻的位置是 $s(t)=t^2$（单位：米，时间：秒）。我们想求 $t=3$ 秒时的瞬时速度。

\begin{itemize}
    \item \textbf{平均速度}：从 $t=3$ 到 $t=3+\Delta t$，位置变化了 $\Delta s = (3+\Delta t)^2 - 3^2 = 6\Delta t + \Delta t^2$，平均速度为 $\frac{\Delta s}{\Delta t} = 6 + \Delta t$。
    \item \textbf{让 $\Delta t$ 越来越小}：比如 $\Delta t=0.1$，平均速度 $6.1$；$\Delta t=0.01$，平均速度 $6.01$；$\Delta t=0.001$，平均速度 $6.001$……我们发现，当 $\Delta t$ 趋于 0 时，平均速度趋于 $6$。这个 $6$ 就是 $t=3$ 时的瞬时速度。
\end{itemize}

\textbf{瞬时速度真的存在吗？} 虽然我们无法直接测量“一瞬间”的速度，但可以用越来越短的时间内的平均速度来无限逼近它，就像用越来越精确的测量逼近真实值。数学上，这个“逼近”的过程就是极限。

这个“趋于”的过程就是极限：
\[
v(3) = \lim_{\Delta t \to 0} \frac{(3+\Delta t)^2 - 3^2}{\Delta t} = \lim_{\Delta t \to 0} (6 + \Delta t) = 6
\]

\subsection*{2.2 导数的定义}

一般地，对于函数 $y=f(x)$，它在 $x$ 处的导数（瞬时变化率）定义为：
\[
f'(x) = \lim_{\Delta x \to 0} \frac{f(x+\Delta x) - f(x)}{\Delta x}
\]
这个式子读作“f prime of x”。

\textbf{直观理解：为什么是比值？}

想象两个场景：
\begin{itemize}
    \item 场景 A：一辆车在 1 秒内走了 100 米。
    \item 场景 B：一辆车在 0.001 秒内走了 0.1 米。
\end{itemize}
只看变化量（100 米 vs 0.1 米），你会觉得 A 快；但算一下比值，A 是 100 米/秒，B 是 100 米/秒——其实一样快！

“率”的本质就是两个变化量的比值。导数要描述“瞬间快慢”，就必须是 $\frac{\Delta y}{\Delta x}$，而不是单纯的 $\Delta y$。

\subsection*{2.3 几何意义}

函数图像上，连接两点 $(x, f(x))$ 和 $(x+\Delta x, f(x+\Delta x))$ 的直线称为\textbf{割线}，它的斜率就是平均变化率。当 $\Delta x \to 0$ 时，割线趋近于\textbf{切线}，导数值就是切线的斜率。所以导数也告诉我们曲线在某点的倾斜程度。

\section*{第三部分：用定义求几个基本函数的导数}

在这一部分，我们完全从导数的定义出发，一步步求出几个最基本函数的导数。这不仅是为了得到公式，更是为了让你亲身体验“极限”如何变成“导数”。

\subsection*{3.1 常数函数 $f(x) = c$（例如 $f(x)=2$）}
\[
\frac{f(x+\Delta x)-f(x)}{\Delta x} = \frac{c - c}{\Delta x} = 0
\]
所以 $f'(x) = \lim_{\Delta x \to 0} 0 = 0$。

\begin{quote}
\textbf{结论}：常数的导数为 0。
\end{quote}

\subsection*{3.2 正比例函数 $f(x) = ax$（$a$ 为常数，例如 $f(x)=3x$）}
\[
\frac{a(x+\Delta x) - a x}{\Delta x} = \frac{a\Delta x}{\Delta x} = a
\]
所以 $f'(x) = \lim_{\Delta x \to 0} a = a$。

\begin{quote}
\textbf{结论}：$ax$ 的导数为 $a$。特别地，当 $a=1$ 时，$x$ 的导数为 1。
\end{quote}

\subsection*{3.3 平方函数 $f(x) = x^2$}
\[
\begin{aligned}
\frac{(x+\Delta x)^2 - x^2}{\Delta x} &= \frac{x^2 + 2x\Delta x + (\Delta x)^2 - x^2}{\Delta x} \\
&= 2x + \Delta x
\end{aligned}
\]
令 $\Delta x \to 0$，高阶小量 $\Delta x$ 消失，得 $f'(x) = 2x$。

\begin{quote}
\textbf{结论}：$x^2$ 的导数为 $2x$。
\end{quote}

\subsection*{3.4 指数函数 $f(x) = e^x$}
这里我们需要先认识一个神奇的常数 $e \approx 2.71828$。它定义为：
\[
e = \lim_{n \to \infty} \left(1 + \frac{1}{n}\right)^n
\]
从这个定义可以推出一个非常重要的极限（先记住它，附录中会有直观解释）：
\[
\lim_{\Delta x \to 0} \frac{e^{\Delta x} - 1}{\Delta x} = 1
\]

现在开始求导：
\[
\begin{aligned}
\frac{e^{x+\Delta x} - e^x}{\Delta x} &= \frac{e^x e^{\Delta x} - e^x}{\Delta x} \\
&= e^x \cdot \frac{e^{\Delta x} - 1}{\Delta x}
\end{aligned}
\]
当 $\Delta x \to 0$ 时，$\frac{e^{\Delta x} - 1}{\Delta x} \to 1$，所以：
\[
(e^x)' = e^x
\]

\begin{quote}
\textbf{结论}：$e^x$ 的导数等于它本身！这意味着这个函数的增长速度永远等于它当前的值，非常奇妙。
\end{quote}

\subsection*{3.5 对数函数 $f(x) = \ln x$}
自然对数 $\ln x$ 是以 $e$ 为底的对数。我们还需要另一个重要极限（同样先记住）：
\[
\lim_{t \to 0} \frac{\ln(1+t)}{t} = 1
\]

求导过程：
\[
\frac{\ln(x+\Delta x) - \ln x}{\Delta x} = \frac{1}{\Delta x} \ln\left(1 + \frac{\Delta x}{x}\right)
\]
令 $t = \frac{\Delta x}{x}$，则当 $\Delta x \to 0$ 时 $t \to 0$，且 $\Delta x = x t$，所以：
\[
\frac{1}{\Delta x} \ln(1+t) = \frac{1}{x t} \ln(1+t) = \frac{1}{x} \cdot \frac{\ln(1+t)}{t}
\]
由重要极限，$\frac{\ln(1+t)}{t} \to 1$，因此：
\[
(\ln x)' = \frac{1}{x}
\]

\begin{quote}
\textbf{结论}：$\ln x$ 的导数为 $\frac{1}{x}$。
\end{quote}

\subsection*{3.6 正弦函数 $f(x) = \sin x$}
这里需要两个关于三角函数的极限（同样来自几何直观，先记住）：
\[
\lim_{\theta \to 0} \frac{\sin\theta}{\theta} = 1, \quad \lim_{\theta \to 0} \frac{1-\cos\theta}{\theta} = 0
\]

利用和角公式 $\sin(x+\Delta x) = \sin x \cos\Delta x + \cos x \sin\Delta x$：
\[
\begin{aligned}
\frac{\sin(x+\Delta x) - \sin x}{\Delta x} &= \frac{\sin x \cos\Delta x + \cos x \sin\Delta x - \sin x}{\Delta x} \\
&= \sin x \cdot \frac{\cos\Delta x - 1}{\Delta x} + \cos x \cdot \frac{\sin\Delta x}{\Delta x}
\end{aligned}
\]
当 $\Delta x \to 0$ 时，$\frac{\cos\Delta x - 1}{\Delta x} \to 0$，$\frac{\sin\Delta x}{\Delta x} \to 1$，所以：
\[
(\sin x)' = \cos x
\]

\begin{quote}
\textbf{结论}：$\sin x$ 的导数为 $\cos x$。
\end{quote}

至此，我们已经通过定义求出了六个基本函数的导数：
\[
\begin{aligned}
(c)' &= 0 \\
(ax)' &= a \\
(x^2)' &= 2x \\
(e^x)' &= e^x \\
(\ln x)' &= \frac{1}{x} \\
(\sin x)' &= \cos x
\end{aligned}
\]
这些是后续所有求导的基石，请务必熟悉它们。

\section*{第四部分：基本求导法则——让求导更高效}

如果每次求导都要回到定义，那会非常麻烦。幸运的是，我们可以从定义推导出一些通用的求导法则，让我们能像搭积木一样组合基本函数的导数。这一部分我们只给出法则和直观理解，严格的推导放在附录中，供有兴趣的同学参考。

\subsection*{4.1 和差法则}
\[
(u \pm v)' = u' \pm v'
\]
\textbf{直观理解}：总的变化率等于各部分变化率之和（或差）。比如两个车队一起跑，总速度就是各自速度之和。

\subsection*{4.2 常数因子法则}
\[
(cu)' = c \cdot u' \quad (c\text{为常数})
\]
\textbf{直观理解}：放大常数倍的变化率，就等于变化率放大同样倍数。

\subsection*{4.3 乘积法则}
\[
(uv)' = u'v + uv'
\]
\textbf{口诀}：前导后不导 + 后导前不导。

\textbf{直观理解}：想象一个矩形，长 $u(x)$，宽 $v(x)$，面积 $S = uv$。当 $x$ 增加一点，长增加 $\Delta u$，宽增加 $\Delta v$，面积的变化 = 新面积 - 旧面积 = $(u+\Delta u)(v+\Delta v) - uv = u\Delta v + v\Delta u + \Delta u\Delta v$。

忽略高阶小量 $\Delta u\Delta v$，面积的变化率 $\frac{\Delta S}{\Delta x} \approx u\frac{\Delta v}{\Delta x} + v\frac{\Delta u}{\Delta x}$。取极限，就得到乘积法则。所以它其实是在说：总面积的变化来自两部分——宽不变时长变化带来的面积，加上长不变时宽变化带来的面积。

\subsection*{4.4 商法则}
\[
\left(\frac{u}{v}\right)' = \frac{u'v - uv'}{v^2} \quad (v \neq 0)
\]
\textbf{口诀}：分母平方，分子上导下不导减上不导下导。

\textbf{直观理解}：可以看作是乘积法则的推论，也可以想象成比例的变化率。

\subsection*{4.5 链式法则（复合函数求导）}
若 $y = f(g(x))$，则：
\[
\frac{dy}{dx} = f'(g(x)) \cdot g'(x)
\]
\textbf{直观理解}：当 $x$ 变化时，内层函数 $u=g(x)$ 以速率 $g'(x)$ 变化，而外层函数 $f$ 随 $u$ 以速率 $f'(u)$ 变化，所以总的速率就是两者相乘，就像“变化的放大倍数”相乘。例如 $y = \sin(2x)$，内层变化率是 2，外层（在对应点）的变化率是 $\cos(2x)$，所以导数为 $2\cos(2x)$。

这些法则是求导的利器，请务必记牢。接下来我们就用它们来扩展我们的导数公式库。

\section*{第五部分：应用法则，扩展导数公式库}

现在，利用上面这些法则和已经求出的基本导数，我们可以轻松地求出更多函数的导数。

\subsection*{5.1 幂函数 $y = x^n$（$n$ 为正整数）}
我们已经知道 $x^2$ 的导数为 $2x$。那么 $x^3$ 呢？可以用乘积法则：
\[
\begin{aligned}
(x^3)' &= (x \cdot x^2)' \\
&= x' \cdot x^2 + x \cdot (x^2)' \\
&= 1 \cdot x^2 + x \cdot 2x \\
&= 3x^2
\end{aligned}
\]

同样，$x^4 = x \cdot x^3$，用乘积法则：
\[
\begin{aligned}
(x^4)' &= 1 \cdot x^3 + x \cdot 3x^2 \\
&= 4x^3
\end{aligned}
\]

看出规律了吗？似乎 $(x^n)' = n x^{n-1}$。我们可以用数学归纳法严格证明这个公式对所有正整数 $n$ 成立：

\begin{enumerate}
    \item \textbf{基础}：$n=1$ 时，$(x^1)' = 1 = 1 \cdot x^{0}$，成立。
    \item \textbf{归纳步骤}：假设 $n=k$ 时成立，即 $(x^k)' = k x^{k-1}$。那么对于 $n=k+1$：
    \[
    \begin{aligned}
    (x^{k+1})' &= (x \cdot x^k)' \\
    &= x' \cdot x^k + x \cdot (x^k)' \\
    &= 1 \cdot x^k + x \cdot k x^{k-1} \\
    &= (k+1)x^k
    \end{aligned}
    \]
    这正是 $n=k+1$ 时的形式。因此公式对一切正整数 $n$ 成立。
\end{enumerate}

\textbf{拓展到负整数}：考虑 $x^{-1} = \frac{1}{x}$，用商法则：
\[
\left(\frac{1}{x}\right)' = \frac{0 \cdot x - 1 \cdot 1}{x^2} = -\frac{1}{x^2} = -1 \cdot x^{-2}
\]
这符合公式 $n x^{n-1}$ 如果令 $n = -1$（因为 $(-1)x^{-2} = -x^{-2}$）。实际上，可以证明幂函数求导公式 $(x^n)' = n x^{n-1}$ 对任意整数 $n$（甚至任意实数）都成立，这里就不深入了。

\subsection*{5.2 一般指数函数 $y = a^x$（$a>0$）}
将 $a^x$ 写成以 $e$ 为底的形式：$a^x = e^{x \ln a}$。这里 $\ln a$ 是一个常数。由链式法则：
\[
\begin{aligned}
(a^x)' &= e^{x \ln a} \cdot \frac{d}{dx}(x \ln a) \\
&= e^{x \ln a} \cdot \ln a \\
&= a^x \ln a
\end{aligned}
\]

\textbf{直观理解}：$\ln a$ 就像是一个“速度调节器”。当 $a = e$ 时，$\ln e = 1$，导数就是自身；当 $a > e$ 时，$\ln a > 1$，增长速度比自身还快；当 $a < e$ 时，$\ln a < 1$，增长速度比自身慢。

\subsection*{5.3 一般对数函数 $y = \log_a x$（$a>0, a\neq1$）}
利用换底公式 $\log_a x = \frac{\ln x}{\ln a}$，$\frac{1}{\ln a}$ 是常数，由常数因子法则：
\[
(\log_a x)' = \frac{1}{\ln a} \cdot (\ln x)' = \frac{1}{\ln a} \cdot \frac{1}{x} = \frac{1}{x \ln a}
\]
特别地，当 $a = e$ 时，$\ln e = 1$，就回到 $(\ln x)' = \frac{1}{x}$。

\subsection*{5.4 余弦函数 $y = \cos x$}
\begin{itemize}
    \item \textbf{方法一}：用定义（类似正弦函数）可以得到 $(\cos x)' = -\sin x$。这里我们略去推导，感兴趣的同学可自行尝试。
    \item \textbf{方法二}：利用关系 $\cos x = \sin\left(x + \frac{\pi}{2}\right)$ 和链式法则。令 $u = x + \frac{\pi}{2}$，则 $\cos x = \sin u$，所以：
    \[
    \begin{aligned}
    (\cos x)' &= \cos u \cdot \frac{du}{dx} \\
    &= \cos\left(x + \frac{\pi}{2}\right) \cdot 1 \\
    &= -\sin x
    \end{aligned}
    \]
    （根据三角诱导公式，$\cos\left(x + \frac{\pi}{2}\right) = -\sin x$）
\end{itemize}

\begin{quote}
\textbf{结论}：$\cos x$ 的导数为 $-\sin x$。
\end{quote}

\subsection*{5.5 其他三角函数（选讲）}
例如正切函数 $\tan x = \frac{\sin x}{\cos x}$，可以用商法则求导：
\[
\begin{aligned}
(\tan x)' &= \frac{(\sin x)' \cos x - \sin x (\cos x)'}{\cos^2 x} \\
&= \frac{\cos x \cdot \cos x - \sin x \cdot (-\sin x)}{\cos^2 x} \\
&= \frac{\cos^2 x + \sin^2 x}{\cos^2 x} \\
&= \frac{1}{\cos^2 x} = \sec^2 x
\end{aligned}
\]
类似地可求 $\cot x, \sec x, \csc x$ 的导数，这里不一一列举。

\subsection*{5.6 复合函数求导练习（链式法则实战）}
链式法则是求导中最重要的技巧之一，我们通过几个例子来熟练掌握它。

\begin{itemize}
    \item \textbf{例 1}：$y = \sin(x^2)$ \\
    外层 $\sin u$，内层 $u = x^2$，则 $y' = \cos(x^2) \cdot 2x = 2x \cos(x^2)$。
    \item \textbf{例 2}：$y = e^{3x}$ \\
    外层 $e^u$，内层 $u = 3x$，则 $y' = e^{3x} \cdot 3 = 3e^{3x}$。
    \item \textbf{例 3}：$y = \ln(2x+1)$ \\
    外层 $\ln u$，内层 $u = 2x+1$，则 $y' = \frac{1}{2x+1} \cdot 2 = \frac{2}{2x+1}$。
    \item \textbf{例 4}：$y = (x^3+1)^5$ \\
    外层 $u^5$，内层 $u = x^3+1$，则 $y' = 5(x^3+1)^4 \cdot 3x^2 = 15x^2 (x^3+1)^4$。
    \item \textbf{例 5}：$y = \cos(3x^2)$ \\
    外层 $\cos u$，内层 $u = 3x^2$，则 $y' = -\sin(3x^2) \cdot 6x = -6x \sin(3x^2)$。
\end{itemize}

\section*{接下来}

恭喜你！你已经掌握了导数的核心内容。接下来，你可以：

\begin{itemize}
    \item \textbf{自由探索} 其他的有趣知识 -> 附录
    \item \textbf{系统梳理} 学到的内容 -> 知识清单
    \item \textbf{尝试运用} 学到的内容 -> 综合练习
\end{itemize}

\end{document}